% ID: SEL-D0072-OA-E01
% ESTADO: validada. S072.
Sumamos $2B$ a ambos miembros. Al calcular sus entradas, $C+2B=I_3$, de modo que $XA=I_3$. Para una fila $(u,v,w)$ de $X$, si la fila correspondiente de $I_3$ es $(d_1,d_2,d_3)$, tenemos
\[
-u+3w=d_1,\qquad 4u-v+w=d_2,\qquad -u+2w=d_3.
\]
Restar la tercera igualdad de la primera da $w=d_1-d_3$. Sustituyendo, $u=2d_1-3d_3$ y $v=9d_1-d_2-13d_3$. Cada fila queda determinada de forma única. Con las tres filas de la identidad resulta
\[
X=A^{-1}=\boxed{\begin{pmatrix}2&9&1\\0&-1&0\\-3&-13&-1\end{pmatrix}}.
\]
